\documentclass[10pt,a4paper]{amsart}
\usepackage[margin=19mm]{geometry}
\usepackage{amsmath,amssymb,mathtools}
\usepackage[scr=rsfs,cal=euler]{mathalfa}
\usepackage{xfrac}
\usepackage[unicode,hidelinks,bookmarks=false]{hyperref}
\newcommand{\ie}{i.e.\ }
\newcommand{\cf}{cf.\ }
\newcommand{\resp}{respectively\ }
\newcommand{\Subsection}{\subsection}
\providecommand{\nobreakdash}{\nobreak\hbox{-}}
\setlength{\emergencystretch}{2em}
\multlinegap0pt
\usepackage{bm}
\usepackage{bbm}
\usepackage{dsfont}
\usepackage{xspace}
\usepackage{cancel}
\usepackage{mathtools}
\usepackage{amsthm}
\makeatletter
\@ifundefined{theorem}{\newtheorem{theorem}{Theorem}}{}
\@ifundefined{assumption}{\newtheorem{assumption}[theorem]{Assumption}}{}
\@ifundefined{definition}{\newtheorem{definition}[theorem]{Definition}}{}
\@ifundefined{lemma}{\newtheorem{lemma}[theorem]{Lemma}}{}
\makeatother
\makeatletter
\expandafter\ifx\csname proof\endcsname\relax
\newenvironment{proof}[1][Proof]{\textbf{#1.} }
\fi
\makeatother
\title[$\ell\_{1}^{2}-\eta\ell\_{2}^{2}$ regularization for sparse re]{$\ell\_{1}^{2}-\eta\ell\_{2}^{2}$ regularization for sparse recovery}
\author{Long Li, Liang Ding}
\date{Source: arXiv:2506.11372v1 (CC BY 4.0)}
\begin{document}
\maketitle

\noindent\emph{Source and licence.} $\ell_{1}^{2}-\eta\ell_{2}^{2}$ regularization for sparse recovery. Version arXiv:2506.11372v1 (CC BY 4.0), distributed under the Creative Commons Attribution 4.0 International licence, as recorded in the arXiv licence metadata for this exact version and shown on the abstract page https://arxiv.org/abs/2506.11372v1. Full rights notice: licenses/arxiv-2506.11372-CC-BY-4.0.txt.\par\smallskip

\noindent\textit{Editorial scope.} Counted unit: the Lemma whose source label is \texttt{lem3.6} in the article (source0.tex lines 547--552, source label \texttt{lem3.6}), delivered together with the article's own complete original proof of it (source0.tex lines 554--605, one \texttt{proof} environment of 52 lines). No mathematics is written, repaired or re-derived: statement and proof are the article's own text line for line. Required context retained uncounted: def2.3 (lines 242--248); ass3.5 (lines 537--543). Every definition, display and label of those lines is retained unchanged. Omitted from this package and not redistributed: 1--241, 249--536, 544--546, 553--553, 606--1589. Nothing inside a retained excerpt is omitted or altered: every retained body is delivered exactly as the article's lines are written, so no omission receipt of any kind accompanies this package. Citations: the retained excerpts carry no citation command at all, so no bibliography entry of the article is transcribed and no reference is redistributed. Reference adaptation: none was needed and none was made; every reference inside the delivered unit resolves inside it, so no unresolved reference or citation is produced. Printed slips kept literal: no printed word, symbol, hypothesis or number of the counted statement or of its proof was altered. Layout and class: the article's own class is \texttt{article} with packages including latexsym, mathrsfs, amsfonts, amsmath, amssymb, indentfirst, subeqnarray, verbatim, graphicx, epstopdf, algorithm, algorithmic, subfigure, color; the wrapper is the lane's pinned \texttt{amsart} editorial wrapper at 10pt on A4, which restates the theorem-like environments the retained text needs and adds the article's own macro definitions that the retained text needs (none), with extra wrapper packages (bm, bbm, dsfont, xspace, cancel, mathtools). The delivered numbering is the unit's own. Assets: no retained span contains a figure environment, a graphics-inclusion command, a PDF-page inclusion, a table or a TikZ picture; no image file is copied into this package, the package declares no asset dependency and no image file is redistributed with it. Licence: version arXiv:2506.11372v1 of this article is distributed under the Creative Commons Attribution 4.0 International licence, as recorded in the arXiv licence metadata for this exact version and shown on the abstract page https://arxiv.org/abs/2506.11372v1. A full rights notice is delivered at licenses/arxiv-2506.11372-CC-BY-4.0.txt. Characters: every retained excerpt is pure ASCII, so the delivered engine, XeLaTeX with its default Latin Modern text font, reports no missing character.
\begin{definition}\label{def2.3}
    We define
    \begin{equation*}
        \mathcal{I}(x^{\dagger})=\left\{i\in\mathbb{N}\ \mid\ x_{i}^{\dagger}\neq 0\right\},
    \end{equation*}
    where $x_{i}^{\dagger}$ is the $i{\rm th}$ component of the $\mathcal{R}_{\eta}$-minimum solution $x^{\dagger}$ of the linear problem $Ax=y$.
\end{definition}

\begin{assumption}\label{ass3.5}
    Let $x^{\dagger}\neq 0$ be a sparse $\mathcal{R}_{\eta}$-minimizing solution to the problem $Ax=y$. We assume that
    \begin{equation*}
        e_{i}\in\mathcal{R}\left(A^{*}\right)\ \forall i\in\mathcal{I}(x^{\dagger}),
    \end{equation*}
    where $e_{i}=(\underbrace{0,\cdots,0,1}_{i},0,\cdots)$ and $\mathcal{I}\left(x^{\dagger}\right)$ is defined in Definition \ref{def2.3}. In other words, for each $i\in\mathcal{I}\left(x^{\dagger}\right)$, there exists an element $\omega_{i}\in\mathcal{D}\left(A^{*}\right)$ such that $e_{i}=A^{*}\omega_{i}$.
\end{assumption}

\begin{lemma}\label{lem3.6}
    Assume that Assumption \ref{ass3.5} holds and that $\mathcal{R}_{\alpha,\beta}\left(x\right)\leq M$ $\left(\alpha>\beta>0\right)$for some constant $M>0$. Then there exist constants $c_{1}>c_{2}$ with $c_{1}>0$ and $c_{3}>0$ such that
    \begin{equation}\label{equ3.4}
        c_{3}\left(\alpha-\beta\right)\left\|x-x^{\dagger}\right\|_{\ell_{1}} \leq\mathcal{R}_{\alpha,\beta}\left(x\right)-\mathcal{R}_{\alpha,\beta}\left(x^{\dagger}\right) +\left(c_{1}\alpha-c_{2}\beta\right)\|Ax-Ax^{\dagger}\|_{Y}
    \end{equation}
\end{lemma}

\begin{proof}
    From the definition of index set $\mathcal{I}\left(x^{\dagger}\right)$, let $\gamma\left(x\right):= \left\|x \right\|_{\ell_{1}}+\left\|x^{\dagger}\right\|_{\ell_{1}}$ and we have
    \begin{equation*}
        \left(\alpha-\beta\right)\left\|x-x^{\dagger}\right\|_{\ell_{1}}\left(\left\|x \right\|_{\ell_{1}}+\left\|x^{\dagger}\right\|_{\ell_{1}}\right)
        =\left(\alpha-\beta\right)\gamma\left(x\right)\left(\sum_{i\in\mathcal{I}\left(x^{\dagger}\right)} \left|x_{i}-x_{i}^{\dagger}\right|^{2}+\sum_{i\notin\mathcal{I}\left(x^{\dagger} \right)}\left|x_{i}\right|^{2}\right).
    \end{equation*}
    Consequently,
    \begin{align}\label{equ3.5}
        &  \quad \left(\alpha-\beta\right)\gamma\left(x\right)\left\|x-x^{\dagger}\right\|_{\ell_{1}}- \left(\mathcal{R}_{\alpha,\beta}\left(x\right)-\mathcal{R}_{\alpha,\beta}\left(x^{\dagger}\right)\right) \nonumber\\
        &  = -\alpha\gamma\left(x\right)\left(\sum_{i\in\mathcal{I}\left(x^{\dagger}\right)}\left|x_{i}\right| -\sum_{i\in\mathcal{I}\left(x^{\dagger}\right)}\left|x_{i}^{\dagger}\right|\right) +\left(\alpha-\beta\right)\gamma\left(x\right)\left(\sum_{i\in\mathcal{I}\left(x^{\dagger}\right)} \left|x_{i}- x_{i}^{\dagger}\right|\right) +\beta\left(T_{1}-T_{2}\right),
    \end{align}
    where
    \begin{align*}
        & T_{1}= \sum_{i}\left|x_{i}\right|^{2} -\sum_{i\notin\mathcal{I}\left(x^{\dagger}\right)}\left|x_{i}\right|^{2} -\sum_{i\in\mathcal{I}\left(x^{\dagger}\right)}\left|x_{i}^{\dagger}\right|^{2} =\sum_{i\in\mathcal{I}\left(x^{\dagger}\right)}\left|x_{i}\right|^{2}- \sum_{i\in\mathcal{I}\left(x^{\dagger}\right)}\left|x_{i}^{\dagger}\right|^{2}, \\
        & T_{2}= \gamma\left(x\right) \left(\sum_{i\notin\mathcal{I}\left(x^{\dagger}\right)}\left|x_{i}\right|\right) -\sum_{i\notin\mathcal{I}\left(x^{\dagger}\right)}\left|x_{i}\right|^{2} =\sum_{i\notin\mathcal{I}\left(x^{\dagger}\right)}\left|x_{i}\right| \left(\gamma\left(x\right)-\left|x_{i} \right|\right).
    \end{align*}
    Note that $T_{2}\geq 0$. Thus, from \eqref{equ3.5}, we derive
    \begin{align}\label{equ3.6}
        \left(\alpha-\beta\right)\gamma\left(x\right)\left\|x-x^{\dagger}\right\|_{\ell_{1}} 
        & \leq \left(\mathcal{R}_{\alpha,\beta}\left(x\right) -\mathcal{R}_{\alpha,\beta}\left(x^{\dagger}\right)\right) +\alpha\gamma\left(x\right)\left(\sum_{i\in\mathcal{I}\left(x^{\dagger}\right)}\left|x_{i}-x_{i}^{\dagger}\right| \right) \nonumber\\
        & \quad +\left(\alpha-\beta\right)\gamma\left(x\right)\left(\sum_{i\in\mathcal{I}\left(x^{\dagger}\right)} \left|x_{i}- x_{i}^{\dagger}\right|\right)+\beta T_{1}.
    \end{align}
    Let $m_{1}$ be a constant upper bound for the terms of the form $\left|x_{i}\right|+ \left|x_{i}^{\dagger}\right|$, and let $m_{2}=\sum_{i}\left|x_{i}^{\dagger}\right|$. Then $0<m_{2}\leq\gamma\left(x\right)$, and we can express
    \begin{equation}\label{equ3.7}
        T_{1}=\sum_{i\in\mathcal{I}\left(x^{\dagger}\right)}\left(\left|x_{i}\right|+ \left|x_{i}^{\dagger}\right|\right)\left(\left|x_{i}\right|- \left|x_{i}^{\dagger}\right|\right) \leq m_{1}\sum_{i\in\mathcal{I}\left(x^{\dagger}\right)}\left|x_{i}- x_{i}^{\dagger}\right|
    \end{equation}
    A combination of \eqref{equ3.6} an \eqref{equ3.7} leads to the following result
    \begin{align}\label{equ3.8}
        \left(\alpha-\beta\right)\left\|x-x^{\dagger}\right\|_{\ell_{1}}
        & \leq \frac{\left(\mathcal{R}_{\alpha,\beta}\left(x\right) -\mathcal{R}_{\alpha,\beta}\left(x^{\dagger}\right)\right)}{\gamma\left(x\right)}+ \left[2\alpha-\left(1-\frac{m_{1}}{\gamma\left(x\right)}\right)\beta\right]\left(\sum_{i\in\mathcal{I} \left(x^{\dagger}\right)}\left|x_{i}-x_{i}^{\dagger}\right|\right) \nonumber\\
        & \leq \frac{\left(\mathcal{R}_{\alpha,\beta}\left(x\right) -\mathcal{R}_{\alpha,\beta}\left(x^{\dagger}\right)\right)}{m_{2}}+ \left[2\alpha-\left(1-\frac{m_{1}}{m_{2}}\right)\beta\right]\left(\sum_{i\in\mathcal{I} \left(x^{\dagger}\right)}\left|x_{i}-x_{i}^{\dagger}\right|\right).
    \end{align}
    Furthermore, based on Assumption \ref{ass3.5}, we have
    \begin{equation*}
        \left|x_{i}-x_{i}^{\dagger}\right|=\left|\left<e_{i},x-x^{\dagger}\right>\right| =\left|\left<\omega_{i},Ax-Ax^{\dagger}\right>\right| \leq \max_{i\in\mathcal{I}\left(x^{\dagger}\right)}\left\|\omega_{i}\right\|_{Y}\left\|Ax-Ax^{\dagger}\right\|_{Y}
    \end{equation*}
    for all $i\in\mathcal{I}\left(x^{\dagger}\right)$. Hence, we obtain
    \begin{equation}\label{equ3.9}
        \sum_{i\in\mathcal{I}\left(x^{\dagger}\right)}\left|x_{i}-x_{i}^{\dagger}\right| \leq \left|\mathcal{I}\left(x^{\dagger}\right)\right| \max_{i\in\mathcal{I}\left(x^{\dagger}\right)}\left\|\omega_{i}\right\|_{Y}\left\|Ax-Ax^{\dagger}\right\|_{Y},
    \end{equation}
    where $\left|\mathcal{I}\left(x^{\dagger}\right)\right|$ denotes the size of the index set $\mathcal{I}\left(x^{\dagger}\right)$. Combining equations \eqref{equ3.8} and\eqref{equ3.9} leads us to the conclusion that
    \begin{align*}
        \left(\alpha-\beta\right)\left\|x-x^{\dagger}\right\|_{\ell_{1}} 
        & \leq \frac{\left(\mathcal{R}_{\alpha,\beta}\left(x\right) -\mathcal{R}_{\alpha,\beta}\left(x^{\dagger}\right)\right)}{m_{2}} \\
        & \quad + \left[2\alpha-\left(1-\frac{m_{1}}{m_{2}}\right)\beta\right]\left|\mathcal{I}\left(x^{\dagger}\right)\right| \max_{i\in\mathcal{I}\left(x^{\dagger}\right)}\left\|\omega_{i}\right\|_{Y}\left\|Ax-Ax^{\dagger}\right\|_{Y},
    \end{align*}
    i.e.,
    \begin{equation*}
        c_{3}\left(\alpha-\beta\right)\left\|x-x^{\dagger}\right\|_{\ell_{1}}\leq \left(\mathcal{R}_{\alpha,\beta}\left(x\right) -\mathcal{R}_{\alpha,\beta}\left(x^{\dagger}\right)\right) + \left(c_{1}\alpha-c_{2}\beta\right)\left\|Ax-Ax^{\dagger}\right\|_{Y},
    \end{equation*}
    where $c_{1}=2m_{2}\left|\mathcal{I}\left(x^{\dagger}\right)\right| \max_{i\in\mathcal{I}\left(x^{\dagger}\right)}\left\|\omega_{i}\right\|_{Y}$, $c_{2}=\left(m_{2}-m_{1}\right)\left|\mathcal{I}\left(x^{\dagger}\right)\right| \max_{i\in\mathcal{I}\left(x^{\dagger}\right)}\left\|\omega_{i}\right\|_{Y}$, and $c_{3}=m_{2}$ with the condition $c_{1}\alpha-c_{2}\beta>0$.
\end{proof}

\end{document}
